% Fragment partage : inclus par liste-contacts.tex (dossier autonome)
% et par dossier-conception-alanya.tex (dossier client assemble).
% Ne pas y mettre de preambule ni de \part : c'est l'appelant qui les pose.

\chapter{Le besoin}

\section{Ce qui existe déjà}

Alanya a aujourd'hui une seule forme d'organisation des contacts~: la liste plate des
\textbf{contacts préférés}, les favoris. Elle est implémentée de bout en bout --- base de
données, API, cache local hors ligne, écran dédié.

Ce qui n'existe pas~: aucun moyen de \emph{regrouper} ces favoris. Un utilisateur qui a
quarante contacts préférés les voit dans une seule liste alphabétique, sans distinction
entre sa famille, ses amis et ses collègues.

\begin{alerte}[Un mot à ne pas confondre]
Dans le code actuel, « groupe » désigne uniquement les \textbf{conversations de groupe}
(\code{conversation.isGroup}). C'est autre chose. Une liste de contacts n'est pas une
conversation~: c'est un rangement personnel, invisible des autres.
\end{alerte}

\section{Ce qu'on ajoute}

Permettre à l'utilisateur de ranger ses contacts préférés dans des \textbf{listes nommées}
--- Famille, Amis, Bureau --- pour deux usages~:

\begin{enumerate}[leftmargin=1.6em]
  \item \textbf{filtrer} l'écran des contacts préférés d'un seul geste~;
  \item \textbf{créer une conversation de groupe} avec toute une liste, sans resélectionner
        les membres un par un.
\end{enumerate}

\section{Les règles produit}

\begin{center}
\begin{tabularx}{\linewidth}{@{}p{4.6cm}X@{}}
\toprule
\thd{Qui peut être membre} & Uniquement des contacts \textbf{déjà en favoris}. Une liste
range ce qui est déjà rangé~; elle n'est pas un second carnet d'adresses. \\
\thd{Appartenance} & \textbf{Plusieurs à plusieurs}. Un contact peut être à la fois en
Famille et en Bureau. \\
\thd{Action groupée} & \textbf{Chat de groupe} uniquement. Pas d'appel de groupe ni de
diffusion de statut à ce stade --- les tables les permettront plus tard sans modification. \\
\thd{Visibilité} & \textbf{Strictement privée}. Personne ne sait dans quelle liste il est
rangé, ni même qu'il l'est. \\
\thd{Emplacement} & Puces de filtre dans l'écran des contacts préférés, \emph{et} un écran
dédié de gestion. \\
\thd{Navigation} & \textbf{Aucun sixième onglet.} La barre est codée en dur à cinq onglets
(\code{lib/screens/home/glass\_nav\_bar.dart}). \\
\bottomrule
\end{tabularx}
\end{center}

\chapter{Les écrans}

\section{Vue d'ensemble}

Trois écrans, un parcours~: on filtre, on gère, on agit.

\begin{center}
  \imageOuCadre{maquettes/listes-contacts.png}{\linewidth}%
    {Les trois écrans : contacts préférés filtrés, gestion des listes, détail d'une liste\\
     \texttt{maquettes/listes-contacts.png}}
\end{center}

\begin{note}[Maquette interactive]
\href{https://claude.ai/code/artifact/d15fb039-5e57-436b-be65-c3579d0dd90c}%
{\texttt{claude.ai/code/artifact/d15fb039}}\\
Source versionnée~: \code{docs/architecture/maquettes/listes-contacts.html}
\end{note}

\section{Écran 1 --- Contacts préférés, filtrés}

L'écran existant gagne une rangée de puces au-dessus de la liste~:

\begin{itemize}[leftmargin=1.4em]
  \item une puce \textbf{« Tous »}, active par défaut~;
  \item une puce \textbf{par liste}, portant sa couleur~;
  \item une puce \textbf{« + Gérer »} qui ouvre l'écran de gestion.
\end{itemize}

La rangée défile horizontalement --- le nombre de listes n'est pas borné. Le filtre se
combine avec la recherche existante~: chercher « Ali » dans le filtre Famille cherche dans
la famille seulement.

Chaque contact affiche les puces des listes auxquelles il appartient. C'est le seul endroit
où l'appartenance multiple se voit d'un coup d'{\oe}il.

\section{Écran 2 --- Gestion des listes}

Une liste par ligne~: pastille de couleur, nom, nombre de membres. Créer, renommer,
recolorer, supprimer. Un bouton flottant « Nouvelle liste ».

Supprimer une liste ne supprime aucun contact --- seulement le rangement.

\section{Écran 3 --- Détail d'une liste}

Les membres de la liste, plus les autres favoris présentés en grisé avec une case à cocher~:
ajouter et retirer se font au même endroit, sans écran intermédiaire.

En bas, un bouton pleine largeur~: \textbf{« Créer un groupe \guillemotleft~Famille~\guillemotright~»}.
Le nom de la liste est pré-rempli et reste modifiable. Un appui ouvre la nouvelle
conversation de groupe.

\begin{note}[Pourquoi la case à cocher plutôt qu'une feuille de sélection]
Montrer les non-membres dans la même liste, en grisé, évite un aller-retour vers un écran
de sélection. L'utilisateur voit d'un seul coup qui est dedans et qui pourrait y être.
\end{note}

\section{Points d'entrée}

\begin{enumerate}[leftmargin=1.6em]
  \item Le profil, entrée « Listes de contacts » à côté de « Contacts préférés »
        (\code{lib/screens/profile/profile\_screen.dart}, méthode
        \code{\_openPreferredContacts}).
  \item La puce « + Gérer » de l'écran 1.
\end{enumerate}
